\documentclass[a4paper,11pt,dvipdfmx]{jsarticle}
\usepackage{amsmath, amssymb, amsthm}
\usepackage{bm}
\usepackage{mathtools}
\usepackage{comment}
\usepackage{algorithm}
\usepackage{algorithmic}
\usepackage{bm}
\usepackage{mathtools}
\usepackage{enumitem}
\mathtoolsset{showonlyrefs}
\usepackage[dvipdfmx, colorlinks=true, linkcolor=blue, citecolor=blue, urlcolor=blue, setpagesize=false]{hyperref}
\usepackage{pxjahyper}
% --- 定理・定義環境 ---
\theoremstyle{definition}
\newtheorem{definition}{定義}[section]
\newtheorem{theorem}[definition]{定理}
\newtheorem{lemma}[definition]{補題}
\newtheorem{assumption}{仮定}
\newtheorem{proposition}[definition]{命題}
\newtheorem{requirement}[definition]{要請}
\newtheorem{example}[definition]{例}
\newcommand{\fU}{\underline{f}}
\newcommand{\gU}{\underline{g}}
\newcommand{\uU}{\underline{u}}
\newcommand{\aU}{\underline{a}}
\newcommand{\bU}{\underline{b}}
\newcommand{\cU}{\underline{c}}
\newcommand{\dU}{\underline{d}}
\newcommand{\vU}{\underline{v}}
\newcommand{\zeroU}{\underline{0}}
\newcommand{\C}[2]{C_{{#1},{#2}}}
\newcommand{\cy}[3]{\mathcal{C}_{{#1},{#2},{#3}}}
\newcommand{\inv}{^{-1}}
\newcommand{\HHX}{\hat{H}_X}
\newcommand{\HHZ}{\hat{H}_Z}
\newcommand{\tr}{^\top}

% =========================================
\begin{document}

\title{}
\author{}
\date{\today}
\maketitle
\tableofcontents
\clearpage
\section{02/04のゼミの内容}
\begin{itemize}
  \item 長さ6のサイクルがあるのではないか：関数$f(x)$と行列$F$の関係について、\cite{kasai2025quantum}とは行と列が逆。、アクティブ部でのみサイクルを考える。
  \item UTCBC以外の長さ8のサイクルを避けるべきか？：今は長さ6以外のサイクルを含まないことを条件にしているが、もしUTCBC以外のサイクルも避ければ、性能が良くなるかも？(ならないかも)、エラーがどのサイクル由来かを調べられるとよい
  \item 整数計画法で求めた特殊解はランダムネスが小さい
  \item APMを使っている時点でランダムネスは減少しているのでAPMに縛らないとランダムネスを増やせる
  \item 他の改良案：Pを小さく(2,3以外の素因数を使う)、ガースを大きく(アクティブ行列の取り出し方を変える。)
\end{itemize}

\section{UTCBCのカウント}
UTCBCを構成する$\rU = [r_0, r_1, r_2, r_3]\in[J]^4, \cU = [c_0, c_1, c_2, c_3]\in[L]^4$は以下の条件を満たす。
\begin{align}
  &c_2 - r_1 \equiv c_0 - r_0\pmod{L/2}\label{02061}\\
  &c_3 - r_2 \equiv c_1 - r_1\pmod{L/2}\label{02062}\\
  &c_3 - r_3 \equiv c_1 - r_0\pmod{L/2}\label{02063}\\
  &r_0\neq r_1, r_1\neq r_2, r_2\neq r_3, r_3\neq r_0\label{02064}\\
  &c_0\neq c_1, c_1\neq c_2, c_2\neq c_3, c_3\neq c_0\label{02065}\\
\end{align}
式\eqref{02062}, \eqref{02063}を整理すると
\begin{align}
  &c_3 - c_1 \equiv r_2 - r_1\pmod{L/2}\\
  &c_3 - c_1 \equiv r_3 - r_0\pmod{L/2}\\
\end{align}
なので、
\begin{align}
  &c_2 - c_0 \equiv r_1 - r_0(=\delta_1\neq 0)\pmod{L/2}\\\
  &c_3 - c_1 \equiv r_2 - r_1 \equiv r_3 - r_0(=\delta_2\neq 0)\pmod{L/2}\label{02066}
\end{align}
が成り立つ。
これを満たす$\rU$は、
\begin{align}
  [0, 1, 2, 1]&((\delta_1, \delta_2) = (1, 1))\\
  [1, 0, 1, 2]&((\delta_1, \delta_2) = (-1, 1))\\
  [1, 2, 1, 0]&((\delta_1, \delta_2) = (1, -1))\\
  [2, 1, 0, 1]&((\delta_1, \delta_2) = (-1, -1))\\
\end{align}
のいずれかである。

つぎに$\cU$の組み合わせを考える。$\rU = [0,1,2,1]$とする。
この時、
\begin{align}
  &c_2 - c_0 \equiv \delta_1\pmod{L/2}\\
  &c_3 - c_1 \equiv \delta_2\pmod{L/2}
\end{align}
であり、$c_0$を決めると、$c_2=c_0+\delta_1, c_0+L/2+\delta_1$の2つ存在する。
次に$c_0, c_2$と異なるように$c_1$を決めると、$c_3=c_1+\delta_2, c_1+L/2+\delta_2$の2つ存在する。

ここで、式\eqref{02065}を考慮する。
\paragraph{$c_2=c_0+\delta_1$かつ$c_3=c_1+\delta_2$のとき}$c_3$が$c_0$または$c_3$と重複してしまう$c_1$が1つ存在する。
\paragraph{$c_2=c_0+\delta_1$かつ$c_3=c_1+L/2+\delta_2$のとき}$c_3$が$c_0$または$c_3$と重複してしまう$c_1$が2つ存在する。
\paragraph{$c_2=c_0+L/2+\delta_1$かつ$c_3=c_1+\delta_2$のとき}$c_3$が$c_0$または$c_3$と重複してしまう$c_1$が2つ存在する。
\paragraph{$c_2=c_0+L/2+\delta_1$かつ$c_3=c_1+L/2+\delta_2$のとき}$c_3$が$c_0$または$c_3$と重複してしまう$c_1$が1つ存在する。

よって、$\cU$の組み合わせは
$12\times 2\times ((10-2)+(10-1))=24\times 17 = 408$通り。
$\rU$は4通りあるので、全部で$4\times 408$通り。
ここでシフトを考えると、同一のサイクルを構成するものが4つずつ存在する(周期2の$\rU, \cU$はない)ので、UTCBCの総数は$4\times 408/4 = 408$通りとなる。

\bibliography{refs}
\bibliographystyle{unsrt} %参考文献出力スタイル
\end{document}